\ifdefined\gkver\else\def\gkver{student}\fi
\documentclass[\gkver]{gaokaozhenti}
\begin{document}

\chapter{2024年全国新课标}
%% paperType: 理综物理部分

\begin{enumerate}
\item
%% number: 1
%% paperName: 2024年全国高考新课标第 1 题 6 分
%% typeId: 1
%% type: 单选题
%% chapter: 直线运动
%% point: 运动图象
%% method: 无
%% score: 6
%% degree: 850
%% duplicateId: 0
%% body:
一质点做直线运动，下列描述其位移 $x$ 或速度 $v$ 随时间 $t$ 变化的图像中，可能正确的是\xzanswer{C}

\fourchoices[answer=C, ispicture=true, h=2.7cm]
{figs/01a.jpg}
{figs/01b.jpg}
{figs/01c.jpg}
{figs/01d.jpg}

%% answer: C
%% memo:
\memoanswer{
AB．物体做直线运动，位移与时间成函数关系，AB 选项中一个时间对应 2 个以上的位移，故不可能，故 AB 错误；

CD．同理 D 选项中一个时间对应 2 个速度，只有 C 选项速度与时间是成函数关系，故 C 正确，D 错误。

故选 C。
}

\item
%% number: 2
%% paperName: 2024年全国高考新课标第 2 题 6 分
%% typeId: 1
%% type: 单选题
%% chapter: 曲线运动
%% point: 平抛运动
%% method: 无
%% score: 6
%% degree: 750
%% duplicateId: 0
%% body:
福建舰是我国自主设计建造的首艘弹射型航空母舰。借助配重小车可以进行弹射测试，测试时配重小车被弹射器从甲板上水平弹出后，落到海面上。调整弹射装置，使小车水平离开甲板时的动能变为调整前的 4 倍。忽略空气阻力，则小车在海面上的落点与其离开甲板处的水平距离为调整前的\xzanswer{C}

\fourchoices[answer=C]
{0.25 倍}
{0.5 倍}
{2 倍}
{4 倍}

%% answer: C
%% memo:
\memoanswer{
小车水平离开甲板时的动能变为调整前的 4 倍，则由 $E_{k} = \frac{1}{2}mv^{2}$ 可知离开甲板时速度变为调整前的 2 倍；小车离开甲板后做平抛运动，从离开甲板到到达海面上时间不变，根据 $x = vt$ 可知小车在海面上的落点与其离开甲板处的水平距离为调整前的 2 倍。

故选 C。
}

\item
%% number: 3
%% paperName: 2024年全国高考新课标第 3 题 6 分
%% typeId: 1
%% type: 单选题
%% chapter: 曲线运动
%% point: 天体运动
%% method: 无
%% score: 6
%% degree: 550
%% duplicateId: 0
%% body:
天文学家发现，在太阳系外的一颗红矮星有两颗行星绕其运行，其中行星 GJ1002c 的轨道近似为圆，轨道半径约为日地距离的 0.07 倍，周期约为 0.06 年，则这颗红矮星的质量约为太阳质量的\xzanswer{B}

\fourchoices[answer=B]
{0.001 倍}
{0.1 倍}
{10 倍}
{1000 倍}

%% answer: B
%% memo:
\memoanswer{
设红矮星质量为 $M_{1}$，行星质量为 $m_{1}$，半径为 $r_{1}$，周期为 $T_{1}$；太阳的质量为 $M_{2}$，地球质量为 $m_{2}$，到太阳距离为 $r_{2}$，周期为 $T_{2}$；根据万有引力定律有

$G\frac{M_{1}m_{1}}{r_{1}^{2}}=m_{1}\frac{4\pi^{2}}{T_{1}^{2}}r_{1}$，$G\frac{M_{2}m_{2}}{r_{2}^{2}}=m_{2}\frac{4\pi^{2}}{T_{2}^{2}}r_{2}$

联立可得 $\frac{M_{1}}{M_{2}}=\left(\frac{r_{1}}{r_{2}}\right)^{3}\cdot\left(\frac{T_{2}}{T_{1}}\right)^{2}$

由于轨道半径约为日地距离的 0.07 倍，周期约为 0.06 年，可得 $\frac{M_{1}}{M_{2}} \approx 0.1$

故选 B。
}

\item
%% number: 4
%% paperName: 2024年全国高考新课标第 4 题 6 分
%% typeId: 1
%% type: 单选题
%% chapter: 光学
%% point: 光子说
%% method: 无
%% score: 6
%% degree: 650
%% duplicateId: 0
%% body:
三位科学家由于在发现和合成量子点方面的突出贡献，荣获了 2023 年诺贝尔化学奖。不同尺寸的量子点会发出不同颜色的光。现有两种量子点分别发出蓝光和红光，下列说法正确的是\xzanswer{A}

\fourchoices[answer=A]
{蓝光光子的能量大于红光光子的能量}
{蓝光光子的动量小于红光光子的动量}
{在玻璃中传播时，蓝光的速度大于红光的速度}
{蓝光在玻璃中传播时的频率小于它在空气中传播时的频率}

%% answer: A
%% memo:
\memoanswer{
AB．红光的频率小于蓝光的频率，红光的波长大于蓝光的波长，根据 $E = h\nu$ 可知蓝光光子的能量大于红光光子的能量；根据 $p = \frac{h}{\lambda }$ 可知蓝光光子的动量大于红光光子的动量，故 A 正确，B 错误；

C．由于红光的折射率小于蓝光，根据 $v = \frac{c}{n}$ 可知在玻璃中传播时，蓝光的速度小于红光的速度，故 C 错误；

D．光从一种介质射入另一种介质中频率不变，故 D 错误。

故选 A。
}

\item
%% number: 5
%% paperName: 2024年全国高考新课标第 5 题 6 分
%% typeId: 1
%% type: 单选题
%% chapter: 电场
%% point: 带电粒子的平衡
%% method: 无
%% score: 6
%% degree: 450
%% duplicateId: 0
%% body:
如图，两根不可伸长的等长绝缘细绳的上端均系在天花板的 $O$ 点上，下端分别系有均带正电荷的小球 $P$、$Q$；小球处在某一方向水平向右的匀强电场中，平衡时两细绳与竖直方向的夹角大小相等。则\xzanswer{B}

\onepicture[w=4.50cm]{figs/05.png}

\fourchoices[answer=B]
{两绳中的张力大小一定相等}
{$P$ 的质量一定大于 $Q$ 的质量}
{$P$ 的电荷量一定小于 $Q$ 的电荷量}
{$P$ 的电荷量一定大于 $Q$ 的电荷量}

%% answer: B
%% memo:
\memoanswer{
由题意可知，设 $Q$ 和 $P$ 两球之间的库仑力为 $F$，绳子的拉力分别为 $T_{1}$、$T_{2}$，质量分别为 $m_{1}$、$m_{2}$，与竖直方向夹角为 $\theta$。对于小球 $Q$ 有

$q_{1}E + T_{1}\sin\theta = F$，$T_{1}\cos\theta = m_{1}g$

对于小球 $P$ 有

$q_{2}E + F = T_{2}\sin\theta$，$T_{2}\cos\theta = m_{2}g$

联立有 $q_{1}E = F - T_{1}\sin\theta > 0$，$q_{2}E = T_{2}\sin\theta - F > 0$

所以可得 $T_{2} > T_{1}$

又因为 $\frac{T_{1}}{T_{2}}=\frac{m_{1}}{m_{2}}$

可知 $m_{2} > m_{1}$，即 $P$ 的质量一定大于 $Q$ 的质量；两小球的电荷量则无法判断。

故选 B。
}

\item
%% number: 6
%% paperName: 2024年全国高考新课标第 6 题 6 分
%% typeId: 2
%% type: 多选题
%% chapter: 机械波
%% point: 波的图象
%% method: 无
%% score: 6
%% degree: 550
%% duplicateId: 0
%% body:
位于坐标原点 $O$ 的波源在 $t = 0$ 时开始振动，振动图像如图所示，所形成的简谐横波沿 $x$ 轴正方向传播。平衡位置在 $x = 3.5 \Um$ 处的质点 $P$ 开始振动时，波源恰好第 2 次处于波谷位置，则\xzanswer{BC}

\onepicture[w=4.00cm]{figs/06.png}

\fourchoices[answer=BC]
{波的周期是 $0.1 \Us$}
{波的振幅是 $0.2 \Um$}
{波的传播速度是 $10 \Ums$}
{平衡位置在 $x = 4.5 \Um$ 处的质点 $Q$ 开始振动时，质点 $P$ 处于波峰位置}

%% answer: BC
%% memo:
\memoanswer{
AB．波的周期和振幅与波源相同，故可知波的周期为 $T = 0.2 \Us$，振幅为 $A = 0.2 \Um$，故 A 错误，B 正确；

C．$P$ 开始振动时，波源第 2 次到达波谷，故可知此时经过的时间为 $t=\frac{3}{4}T+T=0.35 \Us$

故可得波速为 $v=\frac{x_{OP}}{t}=\frac{3.5}{0.35} \Ums=10 \Ums$，故 C 正确；

D．波从 $P$ 传到 $Q$ 点需要的时间为 $t'=\frac{x_{PQ}}{v}=0.1 \Us=\frac{1}{2}T$

故可知质点 $P$ 处于平衡位置，故 D 错误。

故选 BC。
}

\item
%% number: 7
%% paperName: 2024年全国高考新课标第 7 题 6 分
%% typeId: 2
%% type: 多选题
%% chapter: 电磁感应
%% point: 切割磁感线电动势
%% method: 无
%% score: 6
%% degree: 650
%% duplicateId: 0
%% body:
电动汽车制动时可利用车轮转动将其动能转换成电能储存起来。车轮转动时带动磁极绕固定的线圈旋转，在线圈中产生电流。磁极匀速转动的某瞬间，磁场方向恰与线圈平面垂直，如图所示。将两磁极间的磁场视为匀强磁场，则磁极再转过 $90^\circ$ 时，线圈中\xzanswer{BD}

\onepicture[w=5.00cm]{figs/07.png}

\fourchoices[answer=BD]
{电流最小}
{电流最大}
{电流方向由 $P$ 指向 $Q$}
{电流方向由 $Q$ 指向 $P$}

%% answer: BD
%% memo:
\memoanswer{
如图开始线圈处于中性面位置，当磁极再转过 $90^\circ$ 时，此时穿过线圈的磁通量为 0，故可知电流最大；在磁极转动的过程中，穿过线圈的磁通量在减小，根据楞次定律可知，此时感应电流方向由 $Q$ 指向 $P$。

故选 BD。
}

\item
%% number: 8
%% paperName: 2024年全国高考新课标第 8 题 6 分
%% typeId: 2
%% type: 多选题
%% chapter: 气体性质
%% point: 热力学第一定律
%% method: 无
%% score: 6
%% degree: 450
%% duplicateId: 0
%% body:
如图，一定量理想气体的循环由下面 4 个过程组成：$1\to2$ 为绝热过程（过程中气体不与外界交换热量），$2\to3$ 为等压过程，$3\to4$ 为绝热过程，$4\to1$ 为等容过程。上述四个过程是四冲程柴油机工作循环的主要过程。下列说法正确的是\xzanswer{AD}

\onepicture[w=4.00cm]{figs/08.png}

\fourchoices[answer=AD]
{$1\to2$ 过程中，气体内能增加}
{$2\to3$ 过程中，气体向外放热}
{$3\to4$ 过程中，气体内能不变}
{$4\to1$ 过程中，气体向外放热}

%% answer: AD
%% memo:
\memoanswer{
A．$1\to2$ 为绝热过程，根据热力学第一定律 $\Delta U = Q + W$ 可知此时气体体积减小，外界对气体做功，故内能增加，故 A 正确；

B．$2\to3$ 为等压过程，根据盖-吕萨克定律可知气体体积增大时温度增加，内能增大，此时气体体积增大，气体对外界做功 $W < 0$，故气体吸收热量，故 B 错误；

C．$3\to4$ 为绝热过程，此时气体体积增大，气体对外界做功 $W < 0$，根据热力学第一定律可知气体内能减小，故 C 错误；

D．$4\to1$ 为等容过程，根据查理定律可知压强减小时温度减小，故内能减小，由于体积不变 $W = 0$，故可知气体向外放热，故 D 正确。

故选 AD。
}

\item
%% number: 9
%% paperName: 2024年全国高考新课标第 9 题 6 分
%% typeId: 4
%% type: 实验题
%% chapter: 运动和力
%% point: 验证动量守恒实验
%% method: 无
%% score: 6
%% degree: 650
%% duplicateId: 0
%% body:
某同学用如图所示的装置验证动量守恒定律。将斜槽轨道固定在水平桌面上，轨道末段水平，右侧端点在水平木板上的垂直投影为 $O$，木板上叠放着白纸和复写纸。实验时先将小球 $a$ 从斜槽轨道上 $Q$ 处由静止释放，$a$ 从轨道右端水平飞出后落在木板上；重复多次，测出落点的平均位置 $P$ 与 $O$ 点的距离 $x$，将与 $a$ 半径相等的小球 $b$ 置于轨道右侧端点，再将小球 $a$ 从 $Q$ 处由静止释放，两球碰撞后均落在木板上；重复多次，分别测出 $a$、$b$ 两球落点的平均位置 $M$、$N$ 与 $O$ 点的距离 $x_{M}$、$x_{N}$。

完成下列填空：
\begin{enumerate}
	\item
	记 $a$、$b$ 两球的质量分别为 $m_{a}$、$m_{b}$，实验中须满足条件 $m_{a}$\tkanswer{>}$m_{b}$（填“$>$”或“$<$”）；
	\item
	如果测得的 $x_{P}$、$x_{M}$、$x_{N}$、$m_{a}$ 和 $m_{b}$ 在实验误差范围内满足关系式\tkanswer{$m_{a}x_{P} = m_{a}x_{M} + m_{b}x_{N}$}，则验证了两小球在碰撞中满足动量守恒定律。实验中，用小球落点与 $O$ 点的距离来代替小球水平飞出时的速度，依据是\tkanswer{小球离开斜槽末端后做平抛运动，竖直方向高度相同故下落时间相同，水平方向做匀速直线运动，小球水平飞出时的速度与平抛运动的水平位移成正比。}。
\end{enumerate}
\onepicture[w=6.00cm, align=right]{figs/09.png}
%% answer:
\jdanswer{
\begin{enumerate}
	\item
	$>$

	\item
	$m_{a}x_{P} = m_{a}x_{M} + m_{b}x_{N}$，小球离开斜槽末端后做平抛运动，竖直方向高度相同故下落时间相同，水平方向做匀速直线运动，小球水平飞出时的速度与平抛运动的水平位移成正比。

\end{enumerate}
}
%% memo:
\memoanswer{
（1）为了保证小球碰撞为对心正碰，且碰后不反弹，要求 $m_{a} > m_{b}$；

（2）两球离开斜槽后做平抛运动，由于抛出点的高度相等，它们做平抛运动的时间 $t$ 相等，碰撞前 $a$ 球的速度大小 $v_{0} = \frac{x_{P}}{t}$

碰撞后 $a$ 的速度大小 $v_{a} = \frac{x_{M}}{t}$

碰撞后 $b$ 球的速度大小 $v_{b} = \frac{x_{N}}{t}$

如果碰撞过程系统动量守恒，则碰撞前后系统动量相等，则

$m_{a}v_{0} = m_{a}v_{a} + m_{b}v_{b}$

整理得 $m_{a}x_{P} = m_{a}x_{M} + m_{b}x_{N}$

小球离开斜槽末端后做平抛运动，竖直方向高度相同故下落时间相同，水平方向做匀速直线运动，小球水平飞出时的速度与平抛运动的水平位移成正比。
}

\item
%% number: 10
%% paperName: 2024年全国高考新课标第 10 题 12 分
%% typeId: 4
%% type: 实验题
%% chapter: 电路
%% point: 电路实验
%% method: 无
%% score: 12
%% degree: 450
%% duplicateId: 0
%% body:
学生实验小组要测量量程为 $3 \UV$ 的电压表 V 的内阻 $R_{V}$。可选用的器材有：多用电表，电源 $E$（电动势 $5 \UV$），电压表 V$_{1}$（量程 $5 \UV$，内阻约 $3 \UkO$），定值电阻 $R_{0}$（阻值为 $800 \UO$），滑动变阻器 $R_{1}$（最大阻值 $50 \UO$），滑动变阻器 $R_{2}$（最大阻值 $5 \UkO$），开关 S，导线若干。

完成下列填空：
\begin{enumerate}
	\item
	利用多用电表粗测待测电压表的内阻。首先应\tkanswer[1.5cm]{}（把下列实验步骤前的字母按正确操作顺序排列）；
	A．将红、黑表笔短接
	B．调节欧姆调零旋钮，使指针指向零欧姆
	C．将多用电表选择开关置于欧姆挡“$\times$10”位置
	再将多用电表的红、黑表笔分别与待测电压表的\tkanswer[1.5cm]{}（填“正极、负极”或“负极、正极”）相连，欧姆表的指针位置如图~\ref{2024全国新课标10a}中虚线Ⅰ所示。为了减少测量误差，应将选择开关旋转到欧姆挡\tkanswer[1.5cm]{}（填“$\times$1”“$\times$100”或“$\times$1 k”）位置，重新调节后，测量得到指针位置如图~\ref{2024全国新课标10a}中实线Ⅱ所示，则测得的该电压表内阻为\tkanswer[1.5cm]{}$\UkO$（结果保留 1 位小数）；
	\item
	为了提高测量精度，他们设计了如图~\ref{2024全国新课标10b}所示的电路，其中滑动变阻器应选\tkanswer[1.5cm]{}（填“$R_{1}$”或“$R_{2}$”），闭合开关 S 前，滑动变阻器的滑片应置于\tkanswer[1.5cm]{}（填“a”或“b”）端；
	\item
	闭合开关 S，滑动变阻器滑片滑到某一位置时，电压表 V$_{1}$、待测电压表的示数分别为 $U_{1}$、$U$，则待测电压表内阻 $R_{V} = $\tkanswer[1.5cm]{}（用 $U_{1}$、$U$ 和 $R_{0}$ 表示）；
	\item
	测量得到 $U_{1} = 4.20 \UV$，$U = 2.78 \UV$，则待测电压表内阻 $R_{V} = $\tkanswer[1.5cm]{}$\UkO$（结果保留 3 位有效数字）。
\end{enumerate}
\twopicture[labelA=2024全国新课标10a, labelB=2024全国新课标10b, wA=5.40cm, wB=2.90cm, align=right]
{figs/10a.jpg}{figs/10b.jpg}
%% answer:
\jdanswer{
\begin{enumerate}
	\item
	CAB，负极、正极，$\times$100，1.6

	\item
	$R_{1}$，a

	\item
	$\frac{U R_{0}}{U_{1} - U}$

	\item
	1.57

\end{enumerate}
}
%% memo:
\memoanswer{
（1）利用多用电表粗测待测电压表的内阻。首先应选择欧姆挡即 C 选项：将多用电表选择开关置于欧姆挡“$\times$10”位置；接着将红、黑表笔短接即 A 选项；进行欧姆调零即 B 选项：调节欧姆调零旋钮，使指针指向零欧姆。

故首先操作顺序为 CAB。

由多用电表使用时电流“红进黑出”的规则可知：测量电阻时电源在多用电表表内，故将多用电表的红、黑表笔分别与待测电压表的“负极、正极”相连。

读数时欧姆表的指针位置如图（a）中虚线Ⅰ所示，偏转角度较小即倍率选择过小，为了减少测量误差，应将选择开关旋转到欧姆挡倍率较大处，而根据表中数据可知选择“$\times$1 k”倍率又过大，故应选择欧姆挡“$\times$100”的位置；

测量得到指针位置如图（a）中实线Ⅱ所示，则测得的该电压表内阻为 $R = 16.0\times100 \UO = 1.6 \UkO$

（2）图（b）所示的电路，滑动变阻器采用的是分压式连接，为了方便调节，应选最大阻值较小的滑动变阻器即 $R_{1}$；

为保护电路，且测量电路部分电压从零开始，闭合开关 S 前，滑动变阻器的滑片应置于 a 端。

（3）通过待测电压表的电流大小与定值电阻电流相同为 $I = \frac{U_{1} - U}{R_{0}}$

根据欧姆定律得待测电压表的阻值为 $R_{V} = \frac{U}{I} = \frac{U R_{0}}{U_{1} - U}$

（4）测量得到 $U_{1} = 4.20 \UV$，$U = 2.78 \UV$，带入待测电压表的阻值表达式 $R_{V} = \frac{U R_{0}}{U_{1} - U}$，

则待测电压表内阻 $R_{V} = \frac{2.78 \times 800}{4.20 - 2.78} \UO \approx 1.57 \UkO$
}

\item
%% number: 11
%% paperName: 2024年全国高考新课标第 11 题 10 分
%% typeId: 6
%% type: 计算题
%% chapter: 力物体的平衡
%% point: 多力平衡
%% method: 无
%% score: 10
%% degree: 650
%% duplicateId: 0
%% body:
将重物从高层楼房的窗外运到地面时，为安全起见，要求下降过程中重物与楼墙保持一定的距离。如图，一种简单的操作方法是一人在高处控制一端系在重物上的绳子 $P$，另一人在地面控制另一根一端系在重物上的绳子 $Q$，二人配合可使重物缓慢竖直下降。若重物的质量 $m = 42 \Ukg$，重力加速度大小 $g = 10 \Umsq$，当 $P$ 绳与竖直方向的夹角 $\alpha = 37^\circ$ 时，$Q$ 绳与竖直方向的夹角 $\beta = 53^\circ$（$\sin37^\circ = 0.6$）。求：
\begin{enumerate}
	\item
	此时 $P$、$Q$ 绳中拉力的大小；
	\item
	若开始竖直下降时重物距地面的高度 $h = 10 \Um$，在重物下降到地面的过程中，两根绳子拉力对重物做的总功。
\end{enumerate}
\onepicture[w=3.00cm, align=right]{figs/11.png}
%% answer:
\jdanswer{
\begin{enumerate}
	\item
	$1200 \UN$，$900 \UN$

	\item
	$-4200 \UJ$

\end{enumerate}
}
%% memo:
\memoanswer{
（1）重物下降的过程中受力平衡，设此时 $P$、$Q$ 绳中拉力的大小分别为 $T_{1}$ 和 $T_{2}$，竖直方向 $T_{1}\cos\alpha = mg + T_{2}\cos\beta$

水平方向 $T_{1}\sin\alpha = T_{2}\sin\beta$

联立代入数值得 $T_{1} = 1200 \UN$，$T_{2} = 900 \UN$

（2）整个过程根据动能定理得 $W + mgh = 0$

解得两根绳子拉力对重物做的总功为 $W = -4200 \UJ$
}

\item
%% number: 12
%% paperName: 2024年全国高考新课标第 12 题 14 分
%% typeId: 6
%% type: 计算题
%% chapter: 运动和力
%% point: 传送带
%% method: 无
%% score: 14
%% degree: 550
%% duplicateId: 0
%% body:
如图，一长度 $l = 1.0 \Um$ 的均匀薄板初始时静止在一光滑平台上，薄板的右端与平台的边缘 $O$ 对齐。薄板上的一小物块从薄板的左端以某一初速度向右滑动，当薄板运动的距离 $\Delta l = \frac{l}{6}$ 时，物块从薄板右端水平飞出；当物块落到地面时，薄板中心恰好运动到 $O$ 点。已知物块与薄板的质量相等，它们之间的动摩擦因数 $\mu = 0.3$，重力加速度大小 $g = 10 \Umsq$。求：
\begin{enumerate}
	\item
	物块初速度大小及其在薄板上运动的时间；
	\item
	平台距地面的高度。
\end{enumerate}
\onepicture[w=5.00cm, align=right]{figs/12.png}
%% answer:
\jdanswer{
\begin{enumerate}
	\item
	$4 \Ums$，$\frac{1}{3} \Us$

	\item
	$\frac{5}{9} \Um$

\end{enumerate}
}
%% memo:
\memoanswer{
（1）物块在薄板上做匀减速运动的加速度大小为 $a_{1} = \mu g = 3 \Umsq$

薄板做加速运动的加速度 $a_{2} = \frac{\mu mg}{m} = 3 \Umsq$

对物块有 $l + \Delta l = v_{0}t - \frac{1}{2}a_{1}t^{2}$

对薄板有 $\Delta l = \frac{1}{2}a_{2}t^{2}$

解得 $v_{0} = 4 \Ums$，$t = \frac{1}{3} \Us$

（2）物块飞离薄板后薄板的速度为 $v_{2} = a_{2}t = 1 \Ums$

物块飞离薄板后薄板做匀速运动，物块做平抛运动，则当物块落到地面时薄板中心恰好运动到 $O$ 点，所需时间为 $t' = \frac{\frac{l}{2} - \Delta l}{v_{2}} = \frac{\frac{1}{3}l}{v_{2}} = \frac{1}{3} \Us$

则平台距地面的高度 $h = \frac{1}{2}gt'^{2} = \frac{5}{9} \Um$
}

\item
%% number: 13
%% paperName: 2024年全国高考新课标第 13 题 20 分
%% typeId: 6
%% type: 计算题
%% chapter: 磁场
%% point: 带电粒子在组合场中的运动
%% method: 无
%% score: 20
%% degree: 450
%% duplicateId: 0
%% body:
一质量为 $m$、电荷量为 $q$（$q > 0$）的带电粒子始终在同一水平面内运动，其速度可用图示的直角坐标系内一个点 $P(v_{x}, v_{y})$ 表示，$v_{x}$、$v_{y}$ 分别为粒子速度在水平面内两个坐标轴上的分量。粒子出发时 $P$ 位于图中 $a(0, v_{0})$ 点，粒子在水平方向的匀强电场作用下运动，$P$ 点沿线段 $ab$ 移动到 $b(v_{0}, v_{0})$ 点；随后粒子离开电场，进入方向竖直、磁感应强度大小为 $B$ 的匀强磁场，$P$ 点沿以 $O$ 为圆心的圆弧移动至 $c(-v_{0}, v_{0})$ 点；然后粒子离开磁场返回电场，$P$ 点沿线段 $ca$ 回到 $a$ 点。已知任何相等的时间内 $P$ 点沿图中闭合曲线通过的曲线长度都相等。不计重力。求：
\begin{enumerate}
	\item
	粒子在磁场中做圆周运动的半径和周期；
	\item
	电场强度的大小；
	\item
	$P$ 点沿图中闭合曲线移动 1 周回到 $a$ 点时，粒子位移的大小。
\end{enumerate}
\onepicture[w=4.50cm, align=right]{figs/13.png}
%% answer:
\jdanswer{
\begin{enumerate}
	\item
	$\frac{\sqrt{2}mv_{0}}{Bq}$，$\frac{2\pi m}{Bq}$

	\item
	$E = \sqrt{2}Bv_{0}$

	\item
	$\frac{(2 - \sqrt{2})mv_{0}}{Bq}$

\end{enumerate}
}
%% memo:
\memoanswer{
（1）粒子在磁场中做圆周运动时的速度为 $v = \sqrt{v_{0}^{2} + v_{0}^{2}} = \sqrt{2}v_{0}$

根据洛伦兹力提供向心力 $Bqv = m\frac{v^{2}}{r} = m\frac{2\pi}{T}v$

解得做圆周运动的半径为 $r = \frac{\sqrt{2}mv_{0}}{Bq}$，周期为 $T = \frac{2\pi m}{Bq}$

（2）根据题意，已知任何相等的时间内 $P$ 点沿图中闭合曲线通过的曲线长度都相等，由于该曲线表示的为速度相应的曲线，根据 $a = \frac{\Delta v}{\Delta t}$ 可知任意点的加速度大小相等，故可得 $\frac{Bq\cdot\sqrt{2}v_{0}}{m} = \frac{Eq}{m}$

解得 $E = \sqrt{2}Bv_{0}$

（3）根据题意分析可知，从 $b$ 点到 $c$ 点粒子在磁场中转过的角度为 $270^\circ$，绕一圈的过程中两次在电场中运动，根据对称性可知粒子的运动轨迹如图。从 $a$ 到 $b$ 过程中粒子做类平抛运动，得 $\frac{Eq}{m}t = v_{0}$

故可得该段时间内沿 $y$ 方向位移为 $L = v_{0}t$

根据几何知识可得 $bc = 2r$

由粒子在两次电场中运动的对称性可知移动一周时粒子位移的大小为 $x_{aa'} = x_{bc} - 2L$

联立解得 $x_{aa'} = \frac{(2 - \sqrt{2})mv_{0}}{Bq}$
}

\end{enumerate}

\end{document}
